% Chapter 2 -- State of the Art (v3, 2026-10-31: author names only for Ward
% et al.; figures 2.1-2.3 and tables 2.2-2.3 added).
% [VERIFY: ...] marks statements taken from abstracts of papers whose full text
% is not yet in the Literature folder. {[}CITATION{]} marks missing sources.

\label{ch:background}

\tikzset{
  blk/.style={draw, rounded corners=2pt, align=center, inner sep=4pt,
              minimum height=8mm, font=\footnotesize},
  grp/.style={draw, dashed, rounded corners=3pt, inner sep=6pt},
  flow/.style={-{Latex[length=2mm]}, thick},
}

This chapter gives an overview of the current state of research and practice
for recognising assembly activities with sensors worn on the wrist. It starts
with the application itself and with the way industry monitors manual assembly
today. After that, it looks at the two sensor types used in this thesis, the
inertial measurement unit and the microphone, and at the ways their data can be
combined. The last sections deal with running such models on small embedded
devices and with the question of how results in this field are evaluated. The
chapter ends with a comparison of the most closely related work and the
research gap that follows from it.

%------------------------------------------------------------------------------
\section{Activity Recognition in Manual Assembly}
\label{sec:sota-assembly}

Even in highly automated factories, a large part of assembly is still done by
hand~\cite{alamin2019action}. If a system knows which step a worker is
currently performing, it can be used in several ways. It can show new workers
the right instruction at the right time, detect a missing or repeated step for
quality assurance, or measure cycle times without a manual time
study~\cite{stiefmeier2008wearable, alamin2019action}. The same information is
also needed for ergonomic assessments and for human-robot collaboration, where
the robot has to know what its human partner is doing~\cite{berg2018action}.
An overview of these applications in production and logistics is given
in~\cite{reining2019har}.

Assembly activities are different from the activities that most activity
recognition research deals with, such as walking, sitting or climbing stairs.
They are short, often only a few seconds long, they are done mainly with the
hands, and many of them involve a tool. Different activities can look very
similar in terms of arm movement, for example turning a screwdriver and turning
a wrench. In between the activities of interest there are longer phases of
reaching, holding and waiting that do not belong to any class. In the
literature these phases are called the null class~\cite{bulling2014tutorial}.

There are two main ways to recognise such activities, which are shown in
Figure~\ref{fig:sota-sensing}. The first uses cameras that are mounted at the
workstation and record RGB video, depth images or the skeleton of the worker. Deep learning models trained on such data can recognise
the individual steps of an assembly task quite reliably, for example on the
public HA4M dataset of a gear train assembly~\cite{matin2024hybrid,
bonyani2026multimodal}. However, cameras need a clear view of the hands, which
is often blocked by the workpiece or by the worker's own body. They are also
sensitive to viewpoint and lighting~\cite{alamin2019action}, and filming
workers continuously raises privacy concerns. The second way uses sensors that are worn
on the body. They move with the worker, need no free line of sight and measure
the movement directly where it happens. Their disadvantages are that the signal
depends on how the sensor is worn, that the battery is small, and that the data
has to be sent away or processed on a very small device. Large wearable
datasets from industrial settings are still rare. One recent example is
OpenPack, which contains more than 50 hours of packaging work recorded with
wrist-worn and body-worn sensors, depth cameras and barcode
scanners~\cite{yoshimura2024openpack}.

\begin{figure}[!ht]
\centering
\begin{tikzpicture}[
  leaf/.style={blk, text width=2.1cm, minimum height=13mm},
  note/.style={align=left, font=\scriptsize, text width=5.3cm, anchor=north},
]
  \node[blk, text width=4.6cm] (root) at (0,0)
    {Monitoring of manual assembly};
  \node[blk, text width=3.6cm] (stat) at (-3.9,-1.6)
    {\textbf{Stationary}\\installed at the workstation};
  \node[blk, text width=3.6cm] (body) at (2.6,-1.6)
    {\textbf{Body-worn}\\moves with the worker};
  \draw[flow] (root.south) -- ++(0,-0.25) -| (stat.north);
  \draw[flow] (root.south) -- ++(0,-0.25) -| (body.north);

  \node[leaf] (cam)  at (-5.2,-3.5) {\textbf{Cameras}\\RGB, depth,\\skeleton};
  \node[leaf] (tool) at (-2.6,-3.5) {\textbf{Smart tools}\\screwdrivers,\\hand scanners};
  \node[leaf] (imu)  at (0.0,-3.5)  {\textbf{IMU}\\wrist or\\forearm};
  \node[leaf] (mic)  at (2.6,-3.5)  {\textbf{Microphone}\\wrist or\\body};
  \node[leaf] (emg)  at (5.2,-3.5)  {\textbf{EMG}\\forearm\\armband};
  \foreach \n in {cam,tool} \draw[flow] (stat.south) -- ++(0,-0.3) -| (\n.north);
  \foreach \n in {imu,mic,emg} \draw[flow] (body.south) -- ++(0,-0.3) -| (\n.north);

  \node[grp, fit=(imu)(mic), inner sep=4pt,
        label={[font=\scriptsize]below:used in this thesis (on the wrist)}] {};

  \node[note] at (-3.6,-5.0)
    {$+$ no device on the worker\\
     $-$ needs a free view of the hands, privacy concerns\\
     $-$ only connected tools are recorded};
  \node[note] at (2.6,-5.0)
    {$+$ no line of sight needed, works with any tool\\
     $-$ signal depends on how the sensor is worn\\
     $-$ small battery, small computing device};
\end{tikzpicture}
\caption{Approaches for monitoring manual assembly. The dashed box marks the
two sensor types used in this thesis.}
\label{fig:sota-sensing}
\end{figure}

%------------------------------------------------------------------------------
\section{Activity Monitoring in Industrial Practice}
\label{sec:sota-industry}

The idea of using wearable sensors on the assembly line is not new. Already in
2008, a body-worn system for recognising worker activities was tested in two
case studies with the car manufacturer \v{S}koda~\cite{stiefmeier2008wearable}.
One case study dealt with the
training of new assembly workers and the other with quality assurance on the
assembly line. However,
such research prototypes have not led to general activity recognition products.
The systems that are used in factories today record clearly defined events of
special tools and devices instead. Table~\ref{tab:sota-industry} lists some
typical examples.

\begin{table}[!ht]
\centering
\small
\renewcommand{\arraystretch}{1.25}
\begin{tabularx}{\textwidth}{|>{\raggedright\arraybackslash}p{2.9cm}|>{\raggedright\arraybackslash}p{3.4cm}|>{\raggedright\arraybackslash}X|}
\hline
\textbf{Provider} & \textbf{Device} & \textbf{What is recorded} \\
\hline
ProGlove~\cite{proglove2026automotive, proglove2023bmw}
  & Barcode scanner worn on the back of the hand
  & Every scan of a part with a time stamp. Scan and motion data are used to
    analyse process times, bottlenecks and ergonomics. Used at BMW and
    Volkswagen. \\
\hline
Fraunhofer IIS with BMW~\cite{fraunhofer2017tool}
  & Sensor module attached to an electric screwdriver
  & Rotational speed, acceleration and direction of the tool, used to check
    that every screw of a station was tightened. Tested at BMW Regensburg. \\
\hline
Bossard, SAP Digital Manufacturing~\cite{bossard2026smarttool,
  sap2026intelligenttools}
  & Connected screwdrivers and torque tools
  & The tool gets the next work step from the production system and reports
    torque, angle and result of every screw joint. \\
\hline
Ergonomics wearables, e.g.\ LifeBooster~\cite{nsc2025lifebooster}
  & Inertial sensors worn on the upper body and arms
  & Posture and movement over a whole shift, used to estimate the risk of
    musculoskeletal disorders. \\
\hline
\end{tabularx}
\caption[Commercial systems that monitor manual work in production]{Examples
of commercial systems that monitor manual work in production. Information as
given by the providers.}
\label{tab:sota-industry}
\end{table}

A well-known example are the wearable barcode scanners by
ProGlove~\cite{proglove2026automotive, proglove2023bmw}. The scanner
is worn on the back of the hand, so the worker no longer has to pick up and put
down a hand scanner for every part. Car manufacturers such as BMW and
Volkswagen use these scanners in production and logistics, and the related
software analyses scan and motion data to find long process times and
ergonomic problems~\cite{proglove2026automotive, proglove2023bmw}. In the
assembly itself, screwing is monitored most closely, because a missing or
loose screw is a typical quality problem. Fraunhofer IIS developed a sensor
module that can be attached to an existing electric
screwdriver~\cite{fraunhofer2017tool}. From the
rotational speed, acceleration and direction of the tool, the module detects
each screwing operation. In a pilot project at the BMW plant in Regensburg, the
worker was told after each car whether all screws had been
tightened~\cite{fraunhofer2017tool}. Connected screwdrivers, offered for
example by Bossard and integrated into manufacturing software such as SAP
Digital Manufacturing, go one step further. They receive the required torque
from the production system and save the result of every screw
joint~\cite{bossard2026smarttool, sap2026intelligenttools}. A different group
of products uses inertial sensors on the body to assess ergonomic risks over a
whole shift. A report of the U.S.\ Government Accountability Office points out,
however, that the evidence for their benefit comes mostly from the vendors
themselves~\cite{gao2024wearables}.

All of these systems have one thing in common. They either record events that
the tool or device itself produces, such as a scan or a finished screw joint,
or they measure posture without knowing which task is being done. Work with
tools that are not connected, such as a hammer, a manual screwdriver, a file or
a tape measure, is not recorded at all, and every new tool would have to be
equipped and integrated separately. A sensor on the worker's wrist could, in
principle, cover all hand tools at once. This is the starting point for the
research described in the following sections.

%------------------------------------------------------------------------------
\section{Wrist-Worn Inertial Sensing}
\label{sec:sota-imu}

Inertial measurement units (IMUs), which combine a three-axis accelerometer and
a three-axis gyroscope, are the most common sensors for wearable activity
recognition. They are small, cheap and energy-efficient, and they work without
a line of sight~\cite{bulling2014tutorial, kumar2024har}. For assembly work
the wrist is an obvious position. It is close to the hands, it moves with every
tool operation, and most people are used to wearing a watch there. A smartwatch
can tell apart 25 different hand activities, but only if it is worn on the arm
that does the work~\cite{laput2019sensing}. In one of the first industrial
studies, a single IMU was worn on the wrist of workers on an assembly
line~\cite{koskimaki2009activity}. Its two-second windows of acceleration and
angular velocity were classified with a $k$-nearest-neighbour classifier, and a
state machine detected when a work task was finished. The detected start and
end times differed from the true ones by about one
second on average~\cite{koskimaki2009activity}. Other assembly studies use
armbands on the forearm that also measure muscle activity with
electromyography (EMG)~\cite{tao2018worker, alamin2019action}. The forearm is
close to the wrist, but an armband does not follow the rotation of the hand in
the same way, so the two positions are not equivalent.

The basic processing steps for inertial data are well established. The
continuous signal is cut into windows of fixed length, usually one to a few
seconds long and overlapping by 50\,\%, and each window gets one activity
label~\cite{bulling2014tutorial, banos2014window}. Older systems calculated
hand-crafted features for each window, such as mean, variance or frequency
content, and classified them with methods like $k$-nearest neighbours or hidden
Markov models~\cite{koskimaki2009activity, ward2006activity}. Random forests,
which combine the votes of many decision trees, are still a common baseline for
such features~\cite{chatterjee2020laso, koch2024recognition}. Today, neural
networks usually learn the features directly from the raw signal, and they often
do better. In one comparison reported in~\cite{koch2024recognition}, a simple
CNN on smartphone accelerometer data reached 92.7\,\% accuracy, compared with
89.1\,\% for a random forest.
One-dimensional convolutional neural networks (CNNs) move small learned filters
along the time axis and detect typical motion patterns~\cite{tao2018worker,
alamin2019action}. If recurrent layers are added after the convolutions, as in
the DeepConvLSTM architecture, the network can also learn in which order these
patterns occur~\cite{ordonez2016deep, zhou2025efficient}. Attention-based models go further and learn which parts of the signal
are most important~\cite{gao2023mmtsa}. Recent surveys give an overview of
these model types~\cite{kumar2024har}.

For assembly activities that differ clearly in their movement, the reported
results are good. In one study, a CNN recognised six assembly actions, among
them hammering, using a screwdriver and using a wrench, with an accuracy of
98\,\% when part of each worker's data was used for training, and 87\,\% when
the test worker was not part of the training data at
all~\cite{tao2018worker}~[VERIFY: figures against full text]. However, two
problems remain. First, the signal depends on how the device is worn, so a model
trained with one orientation of the device may fail with another
{[}CITATION: sensor orientation and placement in wearable HAR{]}. Second, the
movement of the wrist does not always show what the hand is actually doing.
Turning a screwdriver, tightening a nut with a wrench and turning a knob all
lead to a similar rotation of the forearm. Activities that mainly differ in the
tool or the material are therefore hard to separate by movement alone.
One way to partly solve this is to read the accelerometer of a smartwatch 4000
times per second instead of the usual 50 to 100 times. In this way the sensor
also picks up the small vibrations that a tool or object
transfers to the hand~\cite{laput2019sensing}. This is already a step towards
sound.

%------------------------------------------------------------------------------
\section{Sound as a Cue for Hand and Tool Activity}
\label{sec:sota-audio}

Most tools make a typical sound. A saw makes a regular scraping noise, a hammer
short loud impacts, and a drill or cordless screwdriver the humming of its
motor. A microphone close to the hand records these sounds, together with the
background noise of the workshop. The idea of combining a microphone and an
accelerometer on the arm for this purpose goes back to the
2000s~\cite{lukowicz2004recognizing}. In the best-known study, Ward et~al.\
recognised nine activities of an assembly task in a wood workshop, among them
sawing, hammering and drilling~\cite{ward2006activity}. The user wore
microphones and accelerometers on the arms. The system compared the loudness of
two microphones at different body positions to find the parts of the recording
that contained an activity. The sound was then classified with linear
discriminant analysis, the acceleration with hidden Markov models, and the
results of both were combined. When the system was trained with data of the
same user, it reached a recall of 78\,\% and a precision of 74\,\%. For a new
user these values dropped to 66\,\% and 63\,\%~\cite{ward2006activity}. The
study showed that sound and wrist motion contain different and complementary
information about tool use. It also showed how much harder the task becomes
when the system has to work for people it has not seen before.
Figure~\ref{fig:sota-timeline} shows how this line of research has developed
since then.

\begin{figure}[!ht]
\centering
\begin{tikzpicture}[
  ind/.style={circle, fill, inner sep=2pt},
  daily/.style={circle, draw, thick, fill=white, inner sep=1.7pt},
  mcu/.style={rectangle, fill, inner sep=2.3pt},
]
  \newcommand{\tlrow}[4]{%
    \node[anchor=east, font=\footnotesize] at (0,-#1*0.56) {#2};
    \node[anchor=west, font=\footnotesize, text width=12.6cm]
      at (0.75,-#1*0.56) {#4};
    \node[#3] at (0.4,-#1*0.56) {};}
  \draw[thick] (0.4,0.3) -- (0.4,-3*0.56);
  \draw[thick, dashed] (0.4,-3*0.56) -- (0.4,-4*0.56);
  \draw[thick, -{Latex[length=2mm]}] (0.4,-4*0.56) -- (0.4,-14.6*0.56);
  \tlrow{0}{2004}{ind}{Microphone and accelerometer on the arm for workshop
    activities~\cite{lukowicz2004recognizing}}
  \tlrow{1}{2006}{ind}{Sound and arm motion for an assembly task in a wood
    workshop~\cite{ward2006activity}}
  \tlrow{2}{2008}{ind}{Body-worn activity tracking in car manufacturing at
    \v{S}koda~\cite{stiefmeier2008wearable}}
  \tlrow{3}{2009}{ind}{Single wrist-worn IMU on an industrial assembly
    line~\cite{koskimaki2009activity}}
  \tlrow{4}{2018}{ind}{Forearm IMU and EMG with a CNN for six assembly
    actions~\cite{tao2018worker}}
  \tlrow{5}{2019}{ind}{Fusion levels compared for an assembly step (IMU, EMG,
    camera)~\cite{alamin2019action}}
  \tlrow{6}{2019}{daily}{Smartwatch accelerometer at 4\,kHz picks up hand
    vibrations~\cite{laput2019sensing}}
  \tlrow{7}{2020}{ind}{Workshop audio used to label wrist IMU data
    automatically~\cite{chatterjee2020laso}}
  \tlrow{8}{2022}{daily}{Smartwatch audio and motion, three fusion strategies
    compared~\cite{bhattacharya2022leveraging}}
  \tlrow{9}{2022}{daily}{Smartwatch audio reduced to 1\,kHz, IMU as
    trigger~\cite{mollyn2022samosa}}
  \tlrow{10}{2024}{ind}{OpenPack dataset with wrist IMUs in industrial packaging
    work~\cite{yoshimura2024openpack}}
  \tlrow{11}{2024}{daily}{Ambient audio and vibration in a smart
    home~\cite{koch2024recognition}}
  \tlrow{12}{2024}{mcu}{Camera and IMU fused on a microcontroller
    board~\cite{nakabayashi2024multimodal}}
  \tlrow{13}{2025}{mcu}{Quantised DeepConvLSTM for IMU data on a
    microcontroller~\cite{zhou2025efficient}}
  \tlrow{14}{2026}{ind}{Smartwatch audio and motion for a furniture assembly
    (preprint)~\cite{mahfuz2026proactive}}
  % legend
  \node[ind] at (0.9,-15.7*0.56) {};
  \node[anchor=west, font=\scriptsize] at (1.05,-15.7*0.56)
    {assembly, workshop or industry};
  \node[daily] at (5.6,-15.7*0.56) {};
  \node[anchor=west, font=\scriptsize] at (5.75,-15.7*0.56) {daily life};
  \node[mcu] at (7.9,-15.7*0.56) {};
  \node[anchor=west, font=\scriptsize] at (8.05,-15.7*0.56)
    {focus on microcontrollers};
\end{tikzpicture}
\caption{Development of activity recognition with wrist-worn and body-worn
motion and sound sensors. The dashed part of the time axis marks a gap of
several years without closely related work.}
\label{fig:sota-timeline}
\end{figure}

Since smartwatches contain both an IMU and a microphone, this combination has
received new attention, although mostly for activities of daily life. In one
study, sound and motion were recorded with ordinary smartwatches from 15
participants in their own homes, and 23 activities such as cooking, writing and
cleaning were recognised. Several models and three fusion strategies were
compared, and the combination of both sensors worked better than each sensor
alone. A second study under real everyday conditions showed,
however, that recognition becomes much more difficult outside a controlled
setting~\cite{bhattacharya2022leveraging}~[VERIFY: best fusion strategy in full
text]. The SAMoSA system uses the IMU as a trigger and only
switches on the microphone when movement is detected. In addition, the audio is
recorded with a sampling rate of only 1\,kHz, which saves energy and makes
speech impossible to understand. The system recognised 26 activities in four
indoor environments with an accuracy of 92.2\,\%, and reducing the audio
sampling rate from 16\,kHz to 1\,kHz hardly changed the
result~\cite{mollyn2022samosa}. A recent preprint uses audio and motion from a
smartwatch to follow the steps of a furniture assembly, including screwing and
drilling, and runs the system on edge devices without a cloud
connection~\cite{mahfuz2026proactive}~[VERIFY: sensor setup, preprint].

Sound alone can also be a strong signal. In a smart home study with 25
activities and 14 participants, microphones placed in the rooms reached up to
about 85\,\% accuracy, vibration sensors up to about 80\,\%, and the best
results were achieved with audio or with audio and vibration
together~\cite{koch2024recognition}. As in many audio classification systems,
the sound was first converted into a mel spectrogram, an image-like
representation of how the energy in different frequency bands changes over
time, and then classified with a residual network, a deep
CNN~\cite{koch2024recognition}. On construction sites, combining sound and
motion data of construction machines improved the accuracy by up to 25\,\%
compared with a single sensor~\cite{sherafat2019construction}. Sound can also
help indirectly. In the LASO system, participants wore a smartwatch for motion
data while a smartphone in their pocket recorded audio during workshop
activities such as sawing and hammering. A pretrained audio classifier then
labelled the motion data automatically, and the audio was deleted
afterwards~\cite{chatterjee2020laso}.

The last example shows the main concern about microphones at the workplace.
Audio can contain speech, and most workers will not accept being recorded all
day. In Germany, the introduction of technical devices that can be used to
monitor employees also requires the consent of the works council~\cite{betrvg}. The approaches above deal with this
problem in different ways. Audio is processed on the device and then
deleted~\cite{chatterjee2020laso}, the sampling rate is reduced so far that
speech cannot be understood~\cite{mollyn2022samosa}, or only features or
recognised classes leave the device. Processing the data on the device itself
is therefore not only a technical advantage but also important for the
acceptance of such a system.

%------------------------------------------------------------------------------
\section{Fusion of Inertial and Acoustic Data}
\label{sec:sota-fusion}

IMU and audio data are quite different. A wrist IMU delivers between some tens
and a few hundred samples per second, a microphone several thousand. The IMU
signal is usually used as a time series with a few channels, while audio is
mostly converted into a spectrogram, an image-like representation that shows
how much energy each frequency band contains over time. Before the two signals
can be combined, they also have to be aligned in time, for example by shifting
them until their cross-correlation is highest~\cite{koch2024recognition}. The
most important design decision is at which point in the processing chain the
two signals are combined. Three levels are usually
distinguished~\cite{alamin2019action}.
\begin{itemize}
  \item \textbf{Early fusion} (data-level fusion) combines the raw or
        preprocessed signals into one input for a single model. This requires
        that both signals are brought into a common format.
  \item \textbf{Feature-level fusion} (embedding-level fusion) processes each
        signal with its own encoder network and combines the resulting feature
        vectors, called embeddings, before a common classifier.
  \item \textbf{Late fusion} (decision-level fusion) trains a separate
        classifier for each signal and combines their predicted probabilities,
        for example by taking the average or a weighted sum.
\end{itemize}
In addition, some models use attention mechanisms that learn how much each
signal should contribute. These hybrid approaches do not fit clearly into one
of the three levels~\cite{gao2023mmtsa, nakabayashi2024multimodal}.

Only a few studies compare these levels on the same data with the same sensors.
One of them deals with one step of a 3D printer assembly, recorded with IMU and
EMG armbands on both forearms and a Kinect camera~\cite{alamin2019action}. Five
workers repeated the
step ten times each. A separate CNN was trained for each sensor, and
data-level, feature-level and score-level fusion were compared. A weighted
score fusion, which gives each model more influence on the actions it
recognises best, worked better than a simple average. For fine finger movements
such as tightening a screw, the wearable sensors were more useful than the
camera~\cite{alamin2019action}. In the smart home study with room microphones and
vibration sensors described in Section~\ref{sec:sota-audio}, the data of several
sensor nodes was either merged into one input or processed
separately with the predictions combined afterwards. The separate processing
was clearly better, with up to 85\,\% instead of 75\,\% for audio and 80\,\%
instead of 50\,\% for vibration~\cite{koch2024recognition}. For smartwatch
audio and motion, three fusion strategies were compared
in~\cite{bhattacharya2022leveraging}, and the early work on assembly activities
already tested several ways of combining the decisions of the sound and motion
classifiers~\cite{ward2006activity}~[VERIFY: fusion methods in the full text].

Several studies report that a simple combination does not always make good use
of both signals. In late fusion, one signal can dominate the result while the
other one hardly contributes, an effect called modality
collapse~\cite{bonyani2026multimodal}. Possible solutions are to pretrain the
encoder for each signal separately and fuse the features in several
stages~\cite{choi2022multistage}, or to use contrastive pretraining, where the
model learns that recordings of both signals from the same moment belong
together. The latter also helps when only little labelled data is
available~\cite{ouyang2022cosmo}. Attention can also be used to give the IMU
more weight for activities with body movement and a second sensor more weight
for activities without movement~\cite{nakabayashi2024multimodal}.

The choice of fusion level also affects the computational cost. Early fusion
needs only one network, but both signals have to fit into one common input
format. Feature-level and late fusion keep a separate network for each signal,
so memory and computing time increase. One way to reduce this cost is to
convert the IMU signal into an image, so that the same network can process
both signals~\cite{gao2023mmtsa}. No study was found that compares the fusion
levels for audio and wrist motion and also reports how much memory and time
each variant needs on the same embedded hardware.

%------------------------------------------------------------------------------
\section{Activity Recognition on Microcontrollers}
\label{sec:sota-edge}

The term edge device is used for very different hardware. Single-board
computers such as the Raspberry Pi or the NVIDIA Jetson run a full Linux
system and have several gigabytes of memory, and smartwatches and smartphones
are similar in this respect. Microcontrollers are much smaller. They run
without an operating system or with a small real-time system, have a few
hundred kilobytes to a few megabytes of memory and run at some hundred
megahertz. Compared with smartphones, they have about 100 to 1000 times less
computing power, memory and energy~\cite{david2021tflm}.
Figure~\ref{fig:sota-hardware} shows these hardware classes by their available
memory, together with the devices used in related work. Running neural
networks on such devices is called TinyML. The most widely used framework is
TensorFlow Lite Micro. It executes a converted model using a fixed block of
memory, the so-called tensor arena, which is reserved once at the start, and it
only includes the operations that the model actually
needs~\cite{david2021tflm}. To make the model smaller and to allow integer
arithmetic, the weights and intermediate values are usually converted from
32-bit floating-point numbers to 8-bit integers. This step is called
quantisation. A small set of typical input examples is used to find out the
value range of each layer~\cite{zhou2025efficient}.

\begin{figure}[!ht]
\centering
\begin{tikzpicture}[
  bar/.style={draw, fill=black!12},
  dev/.style={circle, fill, inner sep=1.6pt},
  lab/.style={font=\scriptsize, align=center, inner sep=1pt},
  tier/.style={font=\footnotesize, anchor=east, align=right},
]
  % x = 1.6 cm per decade, x = 0 at 100 KB
  % tiers
  \node[tier] at (-0.25,0)    {Microcontroller};
  \node[tier] at (-0.25,-1.0) {Smartwatch,\\smartphone};
  \node[tier] at (-0.25,-2.0) {Single-board\\computer};
  \node[tier] at (-0.25,-3.0) {Server,\\workstation};
  \draw[bar] (0,-0.17)    rectangle (3.05,0.17);
  \draw[bar] (5.92,-1.17) rectangle (8.13,-0.83);
  \draw[bar] (6.4,-2.17)  rectangle (7.85,-1.83);
  \draw[bar] (8.33,-3.17) rectangle (9.6,-2.83);
  \draw[-{Latex[length=2mm]}, black!50] (9.6,-3) -- (10.1,-3);
  % microcontroller devices
  \node[dev] (nano) at (0.65,0) {};
  \node[lab, anchor=south east] at (0.8,0.2) {Arduino Nano 33\\\cite{zhou2025efficient}};
  \node[dev] (spr) at (2.08,0) {};
  \node[lab, anchor=south west] at (1.95,0.2) {Sony Spresense\\\cite{nakabayashi2024multimodal}};
  \draw[line width=2.2pt] (1.15,0) -- (2.24,0);
  \node[lab, anchor=north, font=\scriptsize\bfseries] at (1.7,-0.22)
    {M5StickC Plus2 (this thesis)};
  % smartwatch
  \node[lab, anchor=south] at (7.0,-0.8)
    {audio and motion studies~\cite{mollyn2022samosa, bhattacharya2022leveraging}};
  % single-board computers
  \node[dev] at (6.4,-2) {};
  \node[lab, anchor=east] at (6.25,-2) {Raspberry Pi 3~\cite{chatterjee2020laso}};
  \node[dev] at (7.85,-2) {};
  \node[lab, anchor=west] at (8.0,-2)
    {Raspberry Pi 4B~\cite{gao2023mmtsa},\\Jetson TX2~\cite{ouyang2022cosmo}};
  % server
  \node[lab, anchor=east] at (8.2,-3) {training and offline evaluation};
  % axis
  \draw[-{Latex[length=2mm]}] (0,-3.7) -- (10.1,-3.7);
  \foreach \x/\t in {0/100\,KB, 1.6/1\,MB, 3.2/10\,MB, 4.8/100\,MB,
                     6.4/1\,GB, 8.0/10\,GB, 9.6/100\,GB} {
    \draw (\x,-3.62) -- (\x,-3.78);
    \node[font=\scriptsize, anchor=north] at (\x,-3.8) {\t};
  }
  \node[font=\footnotesize, anchor=north] at (5.0,-4.15)
    {available working memory (RAM), logarithmic scale};
\end{tikzpicture}
\caption{Hardware classes for running activity recognition models, ordered by
available working memory. Dots mark devices used in related work. The thick
line marks the M5StickC Plus2 used in this thesis, from its 520\,KB of internal
SRAM to an additional 2\,MB of external PSRAM. Memory ranges are typical values
and not exact limits. Memory of the Arduino Nano 33 BLE Sense and the Jetson
TX2 from the manufacturers' data sheets {[}CITATION{]}.}
\label{fig:sota-hardware}
\end{figure}

For models that use only IMU data, this approach is well documented. In one
example, a DeepConvLSTM was trained on the WISDM dataset, which contains
accelerometer data from a smartphone in the trouser pocket, and ran on an
Arduino Nano 33 BLE Sense microcontroller. Quantisation reduced the model size
from 513\,KB to 137\,KB, while the accuracy dropped only from 98.2\,\% to about
97\,\%. On the device, the model needed 29\,KB of RAM and 190\,KB of flash
memory, and one prediction took 21\,ms~\cite{zhou2025efficient}. Multimodal
models are rare on this kind of hardware. One example combines a chest-worn
camera and an IMU with an attention mechanism on a Sony Spresense board with
about 2\,MB of RAM and needs about 21\,ms per prediction. The model size,
quantisation and memory usage are not reported,
however~\cite{nakabayashi2024multimodal}. Many multimodal systems are only
tested on single-board computers such as the Raspberry Pi or the NVIDIA
Jetson~\cite{chatterjee2020laso, gao2023mmtsa, ouyang2022cosmo}.
Table~\ref{tab:sota-deployment} summarises the resource figures that these
studies report. It shows that the figures are often incomplete and that the
models on single-board computers are several orders of magnitude larger than
those on microcontrollers.

\begin{table}[!ht]
\centering
\small
\renewcommand{\arraystretch}{1.2}
\begin{tabularx}{\textwidth}{|>{\raggedright\arraybackslash}p{2.6cm}|>{\raggedright\arraybackslash}p{2.2cm}|>{\raggedright\arraybackslash}X|>{\raggedright\arraybackslash}X|>{\raggedright\arraybackslash}p{2.2cm}|}
\hline
\textbf{Device} & \textbf{Sensors} & \textbf{Model size} & \textbf{Memory used}
  & \textbf{Time per prediction} \\
\hline
Arduino Nano 33 BLE Sense (MCU)~\cite{zhou2025efficient}
  & Accelerometer & 137\,KB, 8-bit integer & 29\,KB RAM, 190\,KB flash & 21\,ms \\
\hline
Sony Spresense (MCU)~\cite{nakabayashi2024multimodal}
  & Camera, IMU & not reported & not reported & about 21\,ms \\
\hline
Raspberry Pi 3 B (SBC)~\cite{chatterjee2020laso}
  & IMU, audio & not reported & CPU and memory load shown, no model size
  & slower than on a laptop$^{a}$ \\
\hline
Raspberry Pi 4B (SBC)~\cite{gao2023mmtsa}
  & Video, accelerometer & 32.4 million parameters, floating point
  & not reported & 5.8\,s for a 15\,s input \\
\hline
NVIDIA Jetson TX2 (SBC)~\cite{ouyang2022cosmo}
  & Depth camera, IMU, radar & not reported & not reported
  & below 0.4\,s for a 2\,s input \\
\hline
\end{tabularx}

\vspace{1mm}
{\footnotesize\raggedright
MCU = microcontroller, SBC = single-board computer.
$^{a}$~Execution time of the whole labelling process, not of a single prediction.\par}
\caption{Resource figures reported for activity recognition on embedded
hardware.}
\label{tab:sota-deployment}
\end{table}

Which type of model is suitable for a microcontroller depends not only on its
accuracy but also on its size and on the operations it needs.
Table~\ref{tab:sota-models} gives an overview of the model types discussed in
this chapter. Hand-crafted features with a random forest need very little
memory, but the features have to be designed by hand for every
sensor~\cite{chatterjee2020laso, koch2024recognition}. CNNs are the most common
compromise. They learn the features themselves, need few parameters because the
same small filters are applied across the whole window, and consist of
operations that TensorFlow Lite Micro supports well~\cite{david2021tflm,
zhou2025efficient}. Recurrent layers such as LSTM or GRU additionally capture
the order of patterns in a window, and they do not have to be large. In a direct
comparison on a microcontroller, a DeepConvLSTM with about 130\,000 parameters
was both more accurate (98.2\,\%) and smaller after quantisation (137\,KB) than
a one-dimensional CNN with about 378\,000 parameters (96.9\,\%,
382\,KB)~\cite{zhou2025efficient}. However, recurrent layers have to process the
time steps one after the other, and TensorFlow Lite Micro supports only a subset
of the operations of the full framework~\cite{david2021tflm}, so recurrent layers
often have to be rewritten into simpler operations before they can run on the
device {[}CITATION{]}. Their benefit therefore has to be weighed against their
cost in each case. Attention-based and transformer models, finally, are mainly
used for multimodal recognition with video~\cite{gao2023mmtsa,
bonyani2026multimodal}. They are much larger. MMTSA, for example, has about 32
million parameters and needs 5.8\,s per input on a Raspberry Pi, which is far
beyond the limits of a microcontroller~\cite{gao2023mmtsa}.

\begin{table}[!ht]
\centering
\small
\renewcommand{\arraystretch}{1.2}
\begin{tabularx}{\textwidth}{|>{\raggedright\arraybackslash}p{3.0cm}|>{\raggedright\arraybackslash}p{2.4cm}|>{\raggedright\arraybackslash}X|>{\raggedright\arraybackslash}p{2.0cm}|}
\hline
\textbf{Model type} & \textbf{Typical input} & \textbf{Suitability for a microcontroller}
  & \textbf{Examples} \\
\hline
Hand-crafted features with a random forest
  & statistics per window, any sensor
  & very small, but features must be designed by hand
  & \cite{chatterjee2020laso, koch2024recognition} \\
\hline
1D CNN
  & raw IMU windows
  & small, simple operations that are well supported
  & \cite{tao2018worker, alamin2019action, zhou2025efficient} \\
\hline
2D CNN on a (mel) spectrogram
  & audio
  & small to medium, but the spectrogram has to be computed on the device
  & \cite{koch2024recognition} \\
\hline
CNN with recurrent layers (e.g.\ DeepConvLSTM)
  & raw IMU windows
  & can be small, but processes time steps one after the other and needs
    operations with limited support
  & \cite{zhou2025efficient} \\
\hline
Attention and transformer models
  & video combined with IMU
  & very large (e.g.\ 32 million parameters~\cite{gao2023mmtsa}), so far only
    on single-board computers or larger
  & \cite{gao2023mmtsa, bonyani2026multimodal} \\
\hline
\end{tabularx}
\caption{Model types for activity recognition and their suitability for
microcontrollers.}
\label{tab:sota-models}
\end{table}

Audio models on microcontrollers are common for keyword spotting, where a
small CNN detects spoken commands in a spectrogram of about one second of audio.
The standard datasets and example applications for this task use a sampling
rate of 16\,kHz~\cite{warden2018speech, zhang2017hello}. If a model works on spectrograms, the feature extraction,
which consists of cutting the signal into short frames, computing a Fourier
transform and applying a mel filter bank, also has to run on the
microcontroller. It has to give the same values as the implementation used for
training, otherwise the model receives inputs it has never seen. For models that
work directly on raw IMU windows, as in the example above, this step is much
simpler~\cite{zhou2025efficient}. When audio and motion are
combined on a wearable, however, the microphone is usually the part that uses
the most energy, which is why SAMoSA only switches it on when
needed~\cite{mollyn2022samosa}. The smartwatch systems for audio and motion
described in Section~\ref{sec:sota-audio} run on smartwatches and phones or are
evaluated offline. No study was found that runs and compares several fusion
variants of audio and wrist motion on a microcontroller.

%------------------------------------------------------------------------------
\section{Evaluation Practice}
\label{sec:sota-evaluation}

How well an activity recognition system seems to work depends strongly on how
the data is split into training and test data. Since neighbouring windows
usually overlap by half, they share a large part of their samples. If windows
are assigned to the training and test set randomly, the test set contains
windows that are almost identical to windows in the training set, and the
measured accuracy is higher than it would be for new recordings. It has been
shown that this effect leads to a clear overestimation of performance for many
different features and classifiers, and that long, repetitive activities are
affected most~\cite{hammerla2015lets}. To avoid it, the data should be split
into continuous blocks instead of single windows.

A second issue is how well a system works for new people. Results are always
higher when data of the same person is used for training and testing than when
the test person is completely left out. Table~\ref{tab:sota-evaluation} shows
this effect for several studies. For the 3D printer assembly, for example, the
best fusion model reached 84.6\,\% when half of the windows of each worker were
used for training, but only 77.8\,\% when the test worker was left out, and the
best single-sensor model reached only about 65\,\% in this
case~\cite{alamin2019action}.

\begin{table}[!ht]
\centering
\small
\renewcommand{\arraystretch}{1.2}
\begin{tabularx}{\textwidth}{|>{\raggedright\arraybackslash}X|>{\raggedright\arraybackslash}p{2.0cm}|>{\raggedright\arraybackslash}p{2.6cm}|>{\raggedright\arraybackslash}p{2.6cm}|}
\hline
\textbf{Task} & \textbf{Metric} & \textbf{Same persons in training and test}
  & \textbf{Test persons left out} \\
\hline
Wood workshop assembly, 9 activities~\cite{ward2006activity}
  & Recall / precision & 78\,\% / 74\,\% & 66\,\% / 63\,\% \\
\hline
Assembly actions with forearm IMU and EMG, 6 actions~\cite{tao2018worker}
  & Accuracy & 98\,\% & 87\,\%$^{a}$ \\
\hline
3D printer assembly step, 7 actions, best fusion model~\cite{alamin2019action}
  & Accuracy & 84.6\,\% & 77.8\,\% \\
\hline
Daily activities (MMAct), video and IMU~\cite{gao2023mmtsa}
  & F1 score & 94.07\,\%$^{b}$ & 87.41\,\% \\
\hline
\end{tabularx}

\vspace{1mm}
{\footnotesize\raggedright
$^{a}$~[VERIFY: figures against full text].
$^{b}$~Split by recording session, so the same persons appear in training and
test.\par}
\caption{Effect of the evaluation protocol on reported results.}
\label{tab:sota-evaluation}
\end{table}

Many studies only give the
ratio of the split, for example 80/20, without saying whether the split was
made by person, by session or by window~\cite{zhou2025efficient,
koch2024recognition, matin2024hybrid}. This makes it hard to compare their
numbers. Datasets from assembly and workshop settings are also small, often with
only about five participants~\cite{ward2006activity, alamin2019action,
chatterjee2020laso}, so the evaluation protocol has a large effect on the
result.
Some methods are designed especially for this situation and are tested with
only one to five percent of the data labelled~\cite{ouyang2022cosmo}.

Finally, assembly datasets are usually imbalanced, because some activities take
much longer than others. The accuracy then mainly shows how well the frequent
classes are recognised. For this reason the macro F1 score is often reported in
addition to accuracy. It is calculated from precision and recall for each class
separately and then averaged, so that every class counts the
same~\cite{bulling2014tutorial}.

%------------------------------------------------------------------------------
\section{Summary and Research Gap}
\label{sec:sota-gap}

Table~\ref{tab:sota-comparison} compares the work that is most closely related
to this thesis. For each system it shows whether the IMU is worn on the wrist,
whether audio is used as an input for classification, whether the activities
involve assembly or tool use, whether different fusion strategies are compared,
and whether the model runs on a microcontroller.

\begin{table}[!ht]
\centering
\small
\renewcommand{\arraystretch}{1.2}
\begin{tabularx}{\textwidth}{|>{\raggedright\arraybackslash}p{4.1cm}|C|C|C|C|C|}
\hline
\textbf{System} & \textbf{Wrist-worn IMU} & \textbf{Audio as input}
  & \textbf{Assembly or tool use} & \textbf{Fusion levels compared}
  & \textbf{Runs on MCU} \\
\hline
Ward et~al.~2006~\cite{ward2006activity}            & yes    & yes    & yes    & partly$^{a}$ & no \\
Stiefmeier et~al.~2008~\cite{stiefmeier2008wearable} & partly$^{b}$ & no & yes  & no     & no \\
Koskim\"aki et~al.~2009~\cite{koskimaki2009activity} & yes    & no     & yes    & --     & no \\
Tao et~al.~2018~\cite{tao2018worker}                & partly$^{c}$ & no & yes  & no     & no \\
Al-Amin et~al.~2019~\cite{alamin2019action}         & partly$^{c}$ & no & yes  & yes    & no \\
Laput and Harrison 2019~\cite{laput2019sensing}     & yes    & partly$^{d}$ & no & --   & no \\
LASO 2020~\cite{chatterjee2020laso}                 & yes    & partly$^{e}$ & partly$^{f}$ & no & no \\
SAMoSA 2022~\cite{mollyn2022samosa}                 & yes    & yes    & no     & no     & no \\
Bhattacharya et~al.~2022~\cite{bhattacharya2022leveraging} & yes & yes & no    & yes    & no \\
OpenPack 2024~\cite{yoshimura2024openpack}          & yes    & no     & partly$^{g}$ & no & no \\
Koch et~al.~2024~\cite{koch2024recognition}         & no     & yes    & no     & partly$^{h}$ & no \\
Nakabayashi and Saito 2024~\cite{nakabayashi2024multimodal} & no & no  & no     & no     & yes \\
Zhou et~al.~2025~\cite{zhou2025efficient}           & no     & no     & no     & --     & yes \\
Mahfuz et~al.~2026~\cite{mahfuz2026proactive}       & yes    & yes    & yes    & no     & no$^{i}$ \\
\hline
This thesis                                         & yes    & yes    & yes    & yes    & yes \\
\hline
\end{tabularx}

\vspace{1mm}
{\footnotesize\raggedright
``--'' not applicable because only one sensor type is used.
$^{a}$~Several ways of combining classifier decisions, no early fusion.
$^{b}$~Several body-worn sensors, not one wrist device.
$^{c}$~Forearm armband with IMU and EMG.
$^{d}$~Accelerometer with a high sampling rate that picks up vibrations, no microphone.
$^{e}$~Audio from a phone in the pocket, only used to label the IMU data.
$^{f}$~Workshop activities such as sawing and hammering, no assembly sequence.
$^{g}$~Packaging work in logistics.
$^{h}$~Merged versus separately processed sensor nodes.
$^{i}$~Runs on edge devices, hardware type not stated in the abstract.\par}
\caption{Comparison of closely related work.}
\label{tab:sota-comparison}
\end{table}

Each of these aspects has been studied on its own. Assembly and tool
activities have been recognised with wrist-worn IMUs since the late
2000s~\cite{ward2006activity, koskimaki2009activity}, and these early
studies already showed that sound recorded close to the hands adds information
that movement alone does not provide. Studies with smartwatches have confirmed that audio and wrist motion
complement each other, and different fusion strategies for these two signals
have been compared. IMU models are regularly quantised and run on
microcontrollers. In industry, on the other hand, assembly is monitored with
connected tools and scanners, and there is no general wearable that recognises
work with any hand tool.

What is missing is a combination of these aspects. The studies that compare
fusion strategies for audio and motion deal with daily activities and run on
smartwatches or offline. The assembly studies that compare fusion levels use
EMG and cameras instead of sound. The multimodal systems that run on a
microcontroller use only one fusion approach and report their resource needs
only in part. It is therefore still open which fusion strategy for audio and
wrist motion is best suited for assembly activities if the model has to run on
the wearable device itself, and how much memory and computing time the
different strategies need on such hardware. This thesis addresses this
question. Chapter~\ref{ch:data} describes the dataset of assembly activities
recorded with wrist-worn devices, and Chapter~\ref{ch:method} describes the
fusion models and their deployment on a microcontroller.
